\FifthArticle{15}
The following remarks on the use of these rules should be added.

To determine each root, two operations must be performed: first, derive from $\lambda$ the corresponding value of $a$ (and hence at the same time that of $b$ or $c$); then calculate from this the value of $x$. For each of these two operations there is a choice of three formulas; in most cases I prefer those listed first. In all these calculations one deals with the logarithms alone, and never with the quantities $\lambda$, $a$, $b$, $c$ themselves. The first operation is indirect and therefore generally involves several successive approximations. It will be convenient to begin with tables having fewer digits. \textsc{Matthiessen}'s table has seven decimal places, as is well known; mine has five; \textsc{Encke} and \textsc{Ursin} have had it printed with four digits. If, at the start of the work, one has no knowledge even of an initial rough approximation, it may be advantageous to have before one, on a separate small sheet, a still shorter extract from the tables with only three digits, for example:
\[
\begin{array}[t]{c|c|}
A&B\\\hline
0&0{,}301\\
0{,}1&0{,}254\\
0{,}2&0{,}212\\
0{,}3&0{,}176\\
0{,}4&0{,}146\\
0{,}5&0{,}119\\
0{,}6&0{,}097\\
0{,}7&0{,}079\\
0{,}8&0{,}064\\
0{,}9&0{,}051\\
1{,}0&0{,}041
\end{array}
\begin{array}[t]{c|c|}
A&B\\\hline
1{,}0&0{,}041\\
1{,}1&0{,}033\\
1{,}2&0{,}027\\
1{,}3&0{,}021\\
1{,}4&0{,}017\\
1{,}5&0{,}014\\
1{,}6&0{,}011\\
1{,}7&0{,}009\\
1{,}8&0{,}007\\
1{,}9&0{,}005\\
2{,}0&0{,}004
\end{array}
\begin{array}[t]{c|c}
A&B\\\hline
2{,}0&0{,}004\\
2{,}1&0{,}003\\
2{,}2&0{,}003\\
2{,}3&0{,}002\\
2{,}4&0{,}002\\
2{,}5&0{,}001\\
2{,}9&0{,}001\\
3{,}0&0{,}000
\end{array}
\]
\SourceLeaf{100}

